\subsection{Tree Models and Bagging}
Tree-based methods partition the input space into simple regions and fit a local rule in each region.  Ensemble averaging methods keep this local modelling idea but stabilise it by resampling, voting, and decorrelating the individual trees.

\begin{definition} \label{def:sl-decision-tree} \hlt{Decision Tree}\\
Let $\mathcal{X}\subseteq\R^p$ and let $\mathcal{A}$ be the prediction space.\\
A decision tree $T$ is a finite binary recursive partition of $\mathcal{X}$ with a leaf value in $\mathcal{A}$ on each terminal region.
\begin{enumerate}[label=(\roman*)]
\setlength{\itemsep}{0pt}
\item \hlt{Split}: a split at node $R\subseteq\mathcal{X}$ is a feature-threshold pair $(j,t)$ with $j\in\{1,\ldots,p\}$ and $t\in\R$.\\
It produces children $R_L\defeq R\cap\{x:x_j\leq t\}$ and $R_R\defeq R\cap\{x:x_j>t\}$.
\item \hlt{Terminal partition}: after splitting stops, the leaves of $T$ are nonempty regions $R_1(T),\ldots,R_M(T)$\\
with $\mathcal{X}=\bigcup\limits_{m=1}^M R_m(T)$ and $R_m(T)\cap R_{m'}(T)=\varnothing$ for $m\neq m'$.
\item \hlt{Leaf rule}: a leaf assignment $a_m\in\mathcal{A}$ for $m=1,\ldots,M$ defines $\widehat{f}_T:\mathcal{X}\to\mathcal{A}$ by $\widehat{f}_T(x)\defeq a_m$\\
for the unique $m$ such that $x\in R_m(T)$.
\end{enumerate}
\end{definition}

\begin{definition} \label{def:sl-bagging-random-forest} \hlt{Bagging and Random Forests}\\
Let $D_n=\{(x_i,y_i)\}_{i=1}^n$, $B\geq1$, and $m\leq p$.
\begin{enumerate}[label=(\roman*)]
\setlength{\itemsep}{0pt}
\item \hlt{Bagging}: draw $D_1^*,\ldots,D_B^*$ from $D_n$, fit $T_1,\ldots,T_B$, and aggregate $\widehat{f}_{\mathrm{bag}}(x)\defeq B^{-1}\sum\limits_{b=1}^B T_b(x)$ for regression, with voting or averaged probabilities for classification.
\item \hlt{Random forest}: in addition to bagging, at each node $R$ of $T_b$ choose $M_{b,R}\subseteq\{1,\ldots,p\}$ with $|M_{b,R}|=m$ and optimise the split over features in $M_{b,R}$ only.
\end{enumerate}
Bagging averages bootstrap perturbations; random split-feature subsets reduce $\operatorname{corr}(T_b,T_{b'})$ before averaging.
\end{definition}

\begin{breakablealgorithm}
\caption{Random Forest}
\label{alg:sl-random-forest}
\begin{algorithmic}
\Require Training data $(x_i,y_i)_{i=1}^n$, number of trees $B$, candidate split count $m\leq p$, tree stopping rule
\Ensure Trees $T_1,\ldots,T_B$ and an aggregated predictor
\For{$b=1,\ldots,B$}
\State Draw a bootstrap resample $S^{\ast b}$ of size $n$
\State Grow a tree $T_b$ on $S^{\ast b}$
\For{each nonterminal node while growing $T_b$}
\State Select $m$ of the $p$ predictors at random
\State Choose the best split using only those $m$ predictors
\State Split the node and continue until the stopping rule is met
\EndFor
\EndFor
\State For regression, return $\widehat{f}_{\mathrm{rf}}(x)=B^{-1}\sum_{b=1}^B T_b(x)$
\State For classification, return a majority vote or averaged class-probability rule from $T_1(x),\ldots,T_B(x)$
\end{algorithmic}
\end{breakablealgorithm}

\begin{figure}[H]
\centering
\begin{tikzpicture}[
    >=Stealth,
    font=\scriptsize,
    line cap=round,
    line join=round,
    box/.style={draw=gray!65, rounded corners=2pt, align=center, minimum height=0.5cm, inner xsep=6pt, inner ysep=2pt},
    data/.style={box, fill=gray!8, minimum width=1.85cm},
    sample/.style={box, draw=cyan!65!black, fill=cyan!8, minimum width=2.05cm},
    tree/.style={box, draw=green!45!black, fill=green!8, minimum width=1.45cm},
    agg/.style={box, draw=orange!70!black, fill=orange!10, minimum width=1.65cm},
    oobbox/.style={box, draw=red!60!black, fill=red!6, minimum width=1.75cm},
    flow/.style={->, draw=black!75, line width=0.6pt},
    oob/.style={->, dashed, draw=red!70!black, line width=0.6pt}
]
\node[data] (data) at (0,0) {training set};
\node[sample] (s1) at (2.55,0.85) {bootstrap $S^{\ast 1}$};
\node[sample] (s2) at (2.55,0) {bootstrap $S^{\ast 2}$};
\node[sample] (s3) at (2.55,-0.85) {bootstrap $S^{\ast 3}$};

\node[tree] (t1) at (5.05,0.85) {tree $T_1$};
\node[tree] (t2) at (5.05,0) {tree $T_2$};
\node[tree] (t3) at (5.05,-0.85) {tree $T_3$};
\node[agg] (agg) at (7.35,0) {aggregate};

\node[oobbox] (held) at (3.15,-1.6) {OOB obs. $i$};
\node[oobbox] (oobavg) at (7.35,-1.6) {OOB prediction};

\draw[flow] (data.north east) -- (s1.west);
\draw[flow] (data.east) -- (s2.west);
\draw[flow] (data.south east) -- (s3.west);
\draw[flow] (s1.east) -- (t1.west);
\draw[flow] (s2.east) -- (t2.west);
\draw[flow] (s3.east) -- (t3.west);
\draw[flow] (t1.east) -- (agg.north west);
\draw[flow] (t2.east) -- (agg.west);
\draw[flow] (t3.east) -- (agg.south west);

\draw[oob] (held.north east) -- (t2.south west);
\draw[oob] (held.east) -- (t3.south west);
\draw[oob] (t2.south east) -- (oobavg.north west);
\draw[oob] (t3.east) -- (oobavg.west);
\end{tikzpicture}

\caption{\textcolor{cyan!65!black}{Blue} bootstraps train \textcolor{green!45!black}{green} trees with $m$ split candidates; black arrows aggregate, \textcolor{red!70!black}{red dashed} arrows give OOB prediction for $i$.}
\label{fig:sl-random-forest-oob-decorrelation}
\end{figure}
